% Transcription — fonds Alexandre Grothendieck, shelfmark 135, batch 3 (pages 41-58).
% Established from the facsimile archives/batches/135/batch-03.pdf (a mirror of
% https://grothendieck.umontpellier.fr/135.pdf), per the skill
% transcribe-grothendieck. Page 45 is blank — absent. Pages 41-44 close the
% annotated typescript begun at page 39 (typed pages 3-6 of Hoàng Xuân Sính's
% "Esquisse d'une théorie des Gr-Catégories"). Pages 46-58 are a thirteen-page
% handwritten manuscript in English, "Gr-categories", paginated 1-13 by its
% author — the inventory's "copies de manuscrit annoté" — in a hand that is not
% the one of the folder's notes, carrying ink and pencil annotations by a
% reader; it ends with a sixteen-entry reference list.
% Pass: Fable 5 (claude-fable-5), 2026-08-15 — first pass, unchecked against the pages by a human.
\documentclass[11pt,a4paper]{article}
% Le préambule commun à toutes les transcriptions du fonds Grothendieck.
%
% Il tient en une idée : **l'incertitude se compose, elle ne se résout pas.**
% Une transcription qui lisse un mot douteux a détruit la seule chose pour
% laquelle on la faisait — un lecteur doit pouvoir savoir, à chaque endroit,
% ce qui était sur le feuillet et ce qui a été ajouté. Les six macros qui
% suivent sont donc l'appareil critique, et non de la décoration.
%
% Les mêmes commandes sont reconnues par `scripts/render.mjs`, qui produit la
% vue de lecture du site. Une macro ajoutée ici doit l'être aussi là-bas,
% faute de quoi le rendu échouera bruyamment — ce qui est voulu : un rendu
% silencieusement faux est pire qu'un rendu absent.

\ProvidesPackage{grothendieck}[2026/08/08 Transcriptions du fonds Grothendieck]

\usepackage{amsmath,amssymb,amsthm}
\usepackage{mathrsfs}
\usepackage{stmaryrd}
% Les diagrammes commutatifs sont partout dans ce fonds : c'est la langue
% même de la géométrie algébrique de Grothendieck, et une transcription qui
% les rendrait en prose ne transcrirait rien.
\usepackage{tikz-cd}
% Français par défaut — c'est la langue des feuillets — mais l'anglais est
% chargé aussi : la traduction et le résumé sont des documents anglais, et la
% typographie française (espaces avant les deux-points) y serait une faute.
% Ces documents basculent par \selectlanguage{english} en tête de corps.
\usepackage[english,french]{babel}
\usepackage{fontspec}
\usepackage{geometry}
\geometry{a4paper,margin=25mm,marginparwidth=28mm}
\usepackage{marginnote}
\usepackage{xcolor}
% ulem, not soul. `soul`'s \st and \ul are built on a character-by-character
% scan that XeLaTeX's font machinery breaks: the document fails with an
% undefined control sequence rather than degrading. ulem does the same two jobs
% and is Unicode-engine safe. `normalem` stops it hijacking \emph, which the
% transcriptions use for Grothendieck's own emphasis.
\usepackage[normalem]{ulem}
\usepackage{hyperref}
\hypersetup{colorlinks=true,linkcolor=black,urlcolor=black,pdfborder={0 0 0}}

% --- Métadonnées du lot -----------------------------------------------------
% Reprises telles quelles par le script de rendu, qui les affiche en tête de
% la vue de lecture. Elles fixent ce que le fichier prétend couvrir : sans
% elles, une transcription flotte sans point d'attache dans un fonds de seize
% mille feuillets.

\newcommand{\folder}[1]{\gdef\grothFolder{#1}}
\newcommand{\batch}[1]{\gdef\grothBatch{#1}}
\newcommand{\leaves}[2]{\gdef\grothFirst{#1}\gdef\grothLast{#2}}
\newcommand{\foldertitle}[1]{\gdef\grothFoldertitle{#1}}
\gdef\grothFolder{?}\gdef\grothBatch{?}\gdef\grothFirst{?}\gdef\grothLast{?}\gdef\grothFoldertitle{}

% --- Le résumé --------------------------------------------------------------
%
% Il ouvre la lecture modernisée, sous le titre « Résumé », et il est composé
% en petit. Ce n'est pas un ornement : c'est le seul passage du document qui
% ne soit pas des mathématiques, et un lecteur doit pouvoir le reconnaître
% avant de l'avoir lu — puis le sauter s'il connaît déjà le sujet, ou s'y
% arrêter s'il ne le connaît pas. Le filet qui le suit marque la reprise du
% corps.

\newenvironment{resume}
  {\small\setlength{\parskip}{0.45em}}
  {\par\medskip\hrule\medskip}

% --- Mots-clés ---------------------------------------------------------------
%
% En anglais, en clôture du résumé de la lecture modernisée : le vocabulaire
% moderne sous lequel chercher ce que les feuillets construisent. Le manifeste
% du site (`npm run manifest`) les extrait pour étiqueter la cote entière —
% cette ligne est la source unique des tags, il n'y a pas de fichier à côté.
\newcommand{\keywords}[1]{\par\smallskip\noindent
  {\footnotesize\textsc{Keywords} --- \textit{#1}}\par}

% --- L'appareil critique ----------------------------------------------------

% Le numéro de feuillet, en marge. C'est l'ancre qui permet de régler un
% désaccord sur une formule en regardant *un* feuillet plutôt qu'une chemise
% de six cents. Le site s'en sert aussi pour tourner les pages du fac-similé
% quand on fait défiler la transcription.
\newcommand{\leaf}[1]{%
  \marginnote{\footnotesize\color{gray!70}\textsf{#1}}%
  \label{leaf:#1}\ignorespaces}

% Illisible : on ne devine pas. Un mot inventé qui se lit comme les autres est
% le pire résultat possible.
\newcommand{\ill}{\textcolor{red!60!black}{[\,\ldots\,]}}

% Lecture douteuse : le mot est donné, et signalé comme incertain.
\newcommand{\uncertain}[1]{\uline{#1}}

% Ajout de l'éditeur — une lettre manquante, un mot sous-entendu.
\newcommand{\add}[1]{\textcolor{blue!50!black}{[#1]}}

% Biffé par Grothendieck lui-même. Ce qu'il a rayé fait partie du document :
% une rature dit souvent où la pensée a changé de direction.
\newcommand{\struck}[1]{\sout{#1}}

% Note du transcripteur, sur l'état matériel du feuillet.
\newcommand{\note}[1]{\par\smallskip{\small\color{gray!60!black}\textit{[#1]}}\par\smallskip}

% Note marginale de Grothendieck, distincte du corps du texte.
\newcommand{\marginal}[1]{\par\smallskip\noindent\hspace{1em}%
  {\small\color{gray!60!black}#1}\par\smallskip}

% --- Titre ------------------------------------------------------------------

\AtBeginDocument{%
  \begin{center}
    {\large\bfseries Fonds Alexandre Grothendieck --- cote n\textsuperscript{o}~\grothFolder}\\[2pt]
    {\itshape\grothFoldertitle}\\[4pt]
    {\small feuillets \grothFirst--\grothLast\ (lot \grothBatch)}\\[2pt]
    {\footnotesize\color{gray!60!black}Transcription établie d'après le fac-similé numérisé
     par l'Université de Montpellier.\\
     Les lectures incertaines sont soulignées, les illisibles notées [\,\ldots\,],
     les ajouts entre crochets.}
  \end{center}
  \vspace{1em}}

\endinput


\folder{135}
\batch{3}
\pages{41}{58}
\dating{[vers 1974]}
\watermark{Édition de démonstration}
\foldertitle{Catégories [et Gr-catégories] : notes manuscrites (s.d.), copies de tapuscrit annoté (s.d.), copies de manuscrit annoté (s.d.).}

\begin{document}

\section*{« Esquisse d'une théorie des Gr-catégories », tapuscrit annoté —
fin (pages 41 à 44)}

\page{41}
\note{la frappe reprend au milieu de la phrase interrompue à la page 40 :
« ... une classe de 3-cohomologie »}
\[
k(C) \in H^{3}(\pi_{0}(C), \pi_{1}(C)).
\]

En précisant ces réflexions, on trouve que la classification, à
$\otimes$-équivalence près, des (Gr)-catégories, se fait précisément en
termes de la donnée d'un groupe $\pi_{0}$, d'un groupe commutatif $\pi_{1}$
sur lequel $\pi_{0}$ opère, et d'une classe de cohomologie (qui peut être
prise arbitraire) dans $H^{3}(\pi_{0}, \pi_{1})$. La loi de \struck{\ill}
groupe du $H^{3}$ admet d'ailleurs une interprétation « géométrique » à la
Baer, en termes d'opérations sur les (Gr)-catégories.

Cas particuliers. Dans l'exemple 1, on doit retrouver bien sûr le premier
invariant de Postnikoff $\in H^{3}(\pi_{1}(X), \pi_{2}(X))$, avec une
interprétation « géométrique » nouvelle. \note{« doit » est un ajout à
l'encre au-dessus de la ligne} Dans l'exemple 2, on trouve $\pi_{0} =
H^{1}(X,F)$, $\pi_{1} = H^{0}(X,F)$ et l'élément du $H^{3}$ est nul (du
fait qu'on a une « catégorie de Picard stricte », cf [Deligne [ ]]).
\note{« Deligne » est écrit à l'encre au-dessus des crochets laissés vides
à la frappe} Dans \struck{\ill} le dernier exemple, on a $\pi_{0}$ = groupe
des classes d'isomorphisme d'autoéquivalences de $A$, $\pi_{1}$ = groupe
des automorphismes du foncteur identique de $A$, et l'invariant $k(C)$
semble nouveau (même dans le cas où $C$ est la catégorie des biModules
inversibles sous un Anneau $O_{X}$, où $\pi_{0}$ peut donc s'interpréter
comme un « groupoïde de Brandt » et $\pi_{1}$ comme le groupe des « unités
centrales » de $O_{X}$). Signalons cependant que lorsque $A$ est la
catégorie des torseurs sous un groupe (ordinaire) $F$, on trouve
\struck{\ill} que $\pi_{0}$ = groupe des automorphismes extérieurs de $F$,
$\pi_{1}$ = centre de $F$, et l'élément $k(C)$ du $H^{3}$ n'est autre que
le classique invariant de Mac-Lane. Il resterait à développer un lien
\uncertain{géométrique} entre la théorie de Mac-Lane des extensions de
groupes non commutatifs, et des « liens » (« kernels » dans sa
terminologie) auxquels ils sont associés, et la théorie \struck{développée
\ill} esquissée ici des (Gr)-catégories et de leurs splittages.
\note{« développer » est écrit à l'encre au-dessus d'un « établir » barré,
et « géométrique » au-dessus de « direct »}

2. Catégories de Picard.

On appelle « catégorie de Picard » une $\otimes$-catégorie ACU
(associative-commutative-unitaire, cf.\ [ ]) qui est une (Gr)-catégorie
pour sa \struck{\ill} $\otimes$-structure AU. On dit que c'est une
catégorie de Picard \emph{stricte} si pour tout $X \in \mathrm{Ob}\,C$, la
symétrie de $X \otimes X$ est l'identité. Ces \struck{\ill} dernières
catégories sont étudiées systématiquement (dans le contexte plus général
des « champs de Picards ») par Deligne [ ].

Soit $C$ une catégorie de Picard. En tant que (Gr)-catégorie, $C$ possède
des invariants $\pi_{0}$ et $\pi_{1}$, mais dans le cas actuel $\pi_{0}$
est aussi commutatif et agit trivialement sur $\pi_{1}$.
\note{l'indice du dernier $\pi$ est surchargé à la frappe}
On fera attention ici qu'une équivalence de (Gr)-catégories entre deux
catégories de Picard $C$, $C'$ n'est pas nécessairement compatible avec
les contraintes de commutativité, donc n'est pas nécessairement une
équivalence de catégories de

\page{42}
Picard — on peut dire que la classification « au sens Picard » est plus
fine que la classification au sens des (Gr)-catégories.

Associant à tout $X \in \mathrm{Ob}\,C$ la symétrie de $X \otimes X$,
interprétée comme \struck{\ill} étant la tensorisation par un élément bien
déterminé $s(X)$ de $\pi_{1}(C) = \mathrm{Aut}(1_{C})$, on trouve un
homomorphisme canonique
\[
s(C) = s : \pi_{0} \longrightarrow {}_{2}(\pi_{1}),
\]
où ${}_{2}(\pi_{1})$ désigne le sous-groupe des éléments $x$ de $\pi_{1}$
tels que $2x = 0$. \note{« homomorphisme » est écrit à l'encre au-dessus
d'un « application » barré} Le résultat principal de la théorie de
structure des catégories de Picard peut s'énoncer en disant que $C$ est
déterminée (à équivalence \uncertain{non canonique} de catégories de
Picard près) par la donnée des groupes commutatifs $\pi_{0}$ et $\pi_{1}$,
et l'homomorphisme $s : \pi_{0} \to {}_{2}(\pi_{1})$. \note{« non
canonique » est un ajout à l'encre} Dire que cet homomorphisme est nul
signifie d'ailleurs que $C$ est une catégorie de Picard stricte, et on
retrouve le résultat de [ ] suivant lequel une catégorie de Picard stricte
est déterminée (à équivalence près) par la connaissance de ses seuls
invariants $\pi_{0}$ et $\pi_{1}$.

Remarques. La théorie de loc.\ cit.\ nous permet, plus précisément,
d'interpréter la 2-catégorie des catégories de Picard strictes en termes
de la sous-catégorie pleine de la catégorie dérivée de (Ab) (cat.\ des
groupes abéliens), formée des complexes $K_{.}$ tels que $H_{i}(K_{.}) =
0$ pour $i \neq 0, 1$. \note{les indices des $H$ sont repassés à l'encre ;
un mot glissé sous la ligne — « alors » ? — reste sans emploi clair} Le
résultat précédent, assez grossier, se réduit à la remarque qu'un objet de
cette catégorie est déterminé à isomorphisme (non unique !) près par la
connaissance de son $H_{0}$ ($= \pi_{0}$) et $H_{1}$ ($= \pi_{1}$). Il y
aurait lieu de développer également pour les catégories de Picard non
nécessairement strictes, et pour les (Gr)-catégories, une théorie de
structure pour la 2-catégorie qu'ils forment. De plus, ces développements
encore dans les limbes, tout comme ceux que \struck{\ill} nous venons
d'esquisser, devraient être un jour \struck{\ill} étudiés dans le cadre
des champs (en (Gr)-catégories, ou en catégories de Picard) tout comme
dans Deligne [ ].
\note{« développements » et « Deligne » sont des ajouts à l'encre, le
premier au-dessus d'un mot barré}

3. Catégories de Picard enveloppantes.

Soient $C$ une \struck{\ill} $\otimes$-catégorie AUC (par exemple une
catégorie avec objet final et admettant des produits cartésiens de deux
éléments), $S$ une partie de $\mathrm{Ob}\,C$. On se pose le problème
universel d'envoyer (par un $\otimes$-foncteur AUC) $C$ dans une
$\otimes$-catégorie AUC $G$, de telle façon que pour tout $X \in S$,
$F(X)$ soit un objet inversible de $G$. \note{« $\otimes$- » et « AUC $G$,
» sont des ajouts à l'encre autour d'un mot barré ; la frappe écrit ici
« AUC » là où la page 41 écrivait « ACU »} Le résultat principal ici est
que ce problème admet toujours une solution, qu'on pourra appeler la
$\otimes$-catégorie AUC « de fractions » par rapport à l'ensemble $S$, et
noter éventuellement $S^{-1}C$. On trouve en fait que pour toute $G$, le
foncteur naturel

\page{43}
\[
\mathrm{Hom}_{\otimes \mathrm{AUC}}(S^{-1}C, G) \longrightarrow
\mathrm{Hom}^{S}_{\otimes \mathrm{AUC}}(C, G)
\]
(où l'exposant $S$ dans le deuxième membre dénote la sous-catégorie pleine
formée des $\otimes$-foncteurs AUC $F : C \to G$ tels que $F(X)$ soit
inversible pour tout $X \in \mathrm{Ob}\,S$) est un isomorphisme de
catégories.

Nous n'explicitons pas ici la construction, assez naturelle, mais dont
l'explicitation pas à pas, avec la vérification de toutes les
compatibilités voulues, est fort longue et fastidieuse. \note{« avec »
remplace, à l'encre, un mot barré} Signalons seulement ici que lorsque $S$
est l'ensemble de tous les objets de $C$, $S^{-1}C$ est elle-même une
catégorie de Picard (pas nécessairement stricte), qu'on pourra appeler la
« catégorie de Picard enveloppante » de $C$. C'est sans doute le cas le
plus important. Les invariants $\pi_{0}$ et $\pi_{1}$ de cette catégorie
enveloppante méritent la notation $K^{0}(C)$ et $K^{1}(C)$ respectivement.
Par exemple, lorsque $C$ est la catégorie des modules projectifs de type
fini sur un anneau $A$, avec l'opération de somme, \struck{on peut
démontrer que} ces groupes ne sont autres que les classiques invariants
$K^{0}(A)$ et $K^{1}(A)$ étudiés par Grothendieck et par Whitehead-Bass.
\note{« avec l'opération de somme, » est un ajout à l'encre, et « on peut
démontrer que » est barré d'un trait ondulé} Le même résultat est sans
doute valable pour toute catégorie additive $C$ (munie de \struck{\ill}
l'opération somme). Cela donne une interprétation remarquable des
invariants $K^{0}$ et $K^{1}$ en termes d'une « catégorie stabilisée » (de
la catégorie additive $C$), laquelle est une catégorie de Picard.
Signalons seulement que, cette catégorie stabilisée (même dans le cas
particulier des modules projectifs de type fini sur l'anneau $A$) n'étant
pratiquement jamais une catégorie de Picard stricte, ces considérations
semblent rendre peu probable qu'on arrive à construire un objet canonique
dans la catégorie dérivée de la catégorie (Ab), qui soit un complexe de
chaînes satisfaisant $H_{0}(K_{.}) = K^{0}$, $H_{1}(K_{.}) = K^{1}$ —
puisqu'une telle construction (grâce à la théorie de Deligne) nous
fournirait précisément une catégorie de Picard stricte canonique dont les
invariants seraient $K^{0}$ et $K^{1}$.

On peut difficilement échapper à la question d'une interprétation
géométrique analogue des invariants $K^{i}$ supérieurs. Il y a tout lieu
de croire qu'une telle interprétation est à chercher dans la théorie des
$n$-catégories de Picard (en un sens convenable) et la construction des
$n$-catégories de Picard enveloppantes d'une $\otimes$-catégorie AUC (ou
d'une $\otimes$-$n$-catégorie AUC) arbitraire $C$, dont les invariants
« homotopiques » $\pi_{i}$ seraient les $K^{i}$ cherchés. De tels
développements obligerai-

\struck{Comme autre direction de recherche ouverte, suggérée par les
résultats esquissés ici, signalons la construction d'une
$\otimes$-catégorie de fractions d'une $\otimes$-catégorie AU (sans donnée
de commutativité néces-}
\note{ce dernier alinéa, barré de grands traits obliques, s'interrompt au
bas de la page sur un mot coupé ; la phrase non barrée qui précède se
poursuit, elle, en tête de la page suivante}

\page{44}
-ent sans doute à tirer au clair les relations très étroites qu'on
pressent entre les notions de $n$-catégorie (et de $\infty$-catégorie)
d'une part, celle d'ensemble simplicial (ou d'espace topologique ...)
d'autre part, tout comme la relation entre $n$-Gr-catégories (et
$\infty$-Gr-catégories) et groupes simpliciaux (ou groupes topologiques),
enfin entre \struck{\ill} $n$-catégories de Picard \emph{strictes} et
\struck{\ill} complexes de chaines tronqués à la dimension $n$. Il y
aurait à s'attendre à une fusion plus ou moins complète entre ces trois
visions : algèbre homologique, algèbre homotopique, algèbre catégorique,
dans une vision commune (qui pourrait bien prendre le nom d'« algèbre
homologique non commutative », étant entendu que ce qui porte actuellement
ce nom n'est qu'une amorce très partielle de ce qui doit venir ...)
\note{« une », « très partielle » et « doit » sont des ajouts à l'encre,
le dernier au-dessus d'un « reste à » barré}

Comme dernière question, plus terre à terre, signalons celle liée à la
remarque \struck{que} suivante (lorsque \struck{\ill}) : $C$ est la
catégorie des modules projectifs de type fini sur un anneau $A$
\emph{commutatif} (ou plus généralement une catégorie additive munie d'une
opération $\otimes$ ACU \emph{biadditive}), alors la catégorie de Picard
enveloppante, en plus de sa loi de composition évidente de « somme »
(notée ici $\oplus$) donnée par sa définition même, est munie d'une
opération $\otimes$ provenant de l'opération de produit tensoriel
ordinaire des modules, qui en fait également une $\otimes$-catégorie ACU
(mais évidemment pas une (Gr)-catégorie), liée par des isomorphismes de
distributivité évidents à l'opération $\oplus$. \note{« suivante »,
« ACU », « par », « ordinaire » et « à l' » sont des retouches à l'encre
sur la frappe} Il y aurait lieu d'étudier systématiquement les catégories
$A$ ainsi munies de deux opérations $\oplus$ et $\otimes$, toutes deux AUC
(des données supplémentaires), avec une donnée de distributivité
\[
X \otimes (Y \oplus Z) \simeq X \otimes Y \oplus X \otimes Z
\]
et des compatibilités convenables pour celles-ci ; et de tenter de faire
une théorie de structure pour de telles catégories qui sont de
\struck{\ill} Picard pour la structure $\oplus$. Cela implique en
particulier que $\pi_{0}(A)$ est un anneau unitaire commutatif, et
$\pi_{1}(A)$ un module sous ce dernier. Quels autres invariants faut-il
introduire pour caractériser $A$ à équivalence près respectant toutes les
structures ?

\section*{« Gr-categories », manuscrit anglais annoté (pages 46 à 58)}

\note{les pages 46 à 58 forment un manuscrit autonome de treize pages,
paginées 1 à 13 par leur auteur, rédigé en anglais : le résumé de
l'organisation d'un travail en trois chapitres sur les Gr-catégories,
s'achevant sur une bibliographie. La main n'est pas celle des notes du
dossier ; des annotations à l'encre et au crayon, signalées au fil du
texte, s'y ajoutent. L'auteur souligne les lettres de catégories et
l'objet unité — soulignement que nous ne reproduisons pas — et nous
transcrivons l'anglais tel quel, sans le traduire}

\page{46}
\textbf{Gr-categories}

\textbf{Summary}

The purpose of these notes is to study the Gr-categories and give some
applications of them. Below is a brief description of the organization of
the work.

Chapter I gives some definitions and results, which are used continually
in the sequel, on $\otimes$-categories one can find in [2], [6], [11],
[14], [15], the terminology employed in this chapter being of Neantro
Saavedra Rivano [14]. A $\otimes$-category is a category $C$ together with
a law $\otimes$, i.e a covariant bifunctor
\[
\otimes : C \times C \longrightarrow C, \qquad
(X, Y) \longmapsto X \otimes Y
\]

An \emph{associativity constraint} for a $\otimes$-category $C$ is an
isomorphism of trifunctors
\[
a_{X,Y,Z} : X \otimes (Y \otimes Z) \xrightarrow{\ \sim\ }
(X \otimes Y) \otimes Z, \qquad X, Y, Z \in \mathrm{Ob}\,C
\]
satisfying the \emph{pentagon axiom}, i.e all the pentagonal diagrams

\begin{tikzcd}[column sep=small, row sep=small, nodes={font=\scriptsize}]
& (X \otimes Y) \otimes (Z \otimes T) \arrow[dr, "a_{X \otimes Y{,}Z{,}T}"]
& \\
X \otimes (Y \otimes (Z \otimes T)) \arrow[ur, "a_{X{,}Y{,}Z \otimes T}"]
\arrow[d, "\mathrm{id}_{X} \otimes a_{Y{,}Z{,}T}"'] & &
((X \otimes Y) \otimes Z) \otimes T \\
X \otimes ((Y \otimes Z) \otimes T) \arrow[rr, "a_{X{,}Y \otimes Z{,}T}"']
& & (X \otimes (Y \otimes Z)) \otimes T
\arrow[u, "a_{X{,}Y{,}Z} \otimes \mathrm{id}_{T}"']
\end{tikzcd}

are commutative. A $\otimes$-category together with an associativity
constraint is called a $\otimes$-\uncertain{associativite} category.
\note{sic — l'orthographe flotte : deux pages plus loin la même plume
écrit « associative »}

A \emph{commutativity constraint} for a $\otimes$-category $C$ is an
isomorphism

\page{47}
of bifunctors
\[
c_{X,Y} : X \otimes Y \xrightarrow{\ \sim\ } Y \otimes X, \qquad
X, Y \in \mathrm{Ob}\,C,
\]
verifying the relation
\[
c_{Y,X} \circ c_{X,Y} = \mathrm{id}_{X \otimes Y}.
\]
The commutativity constraint $c$ is said to be \emph{strict} if $c_{X,X} =
\mathrm{id}_{X \otimes X}$ for all $X \in \mathrm{Ob}\,C$. A
$\otimes$-category together with a commutativity constraint is a
$\otimes$-commutative category. A $\otimes$-commutative category is
\emph{strict} if its commutativity constraint is strict.

An \emph{unity constraint} for a $\otimes$-category $C$ is a triple $(1,
g, d)$ where $1$ is an object of $C$, $g$ and $d$ natural isomorphisms
\[
g_{X} : X \xrightarrow{\ \sim\ } 1 \otimes X, \qquad
d_{X} : X \xrightarrow{\ \sim\ } X \otimes 1, \qquad
X \in \mathrm{Ob}\,C
\]
such that $g_{1} = d_{1}$. A $\otimes$-category together with an unity
constraint is a $\otimes$-unifer category.

A $\otimes$-category $C$ together with an associativity constraint $a$
and a commutativity constraint $c$ is a $\otimes$-AC category if the
\emph{hexagon axiom} is fulfilled, i.e all the hexagonal diagrams commute

\begin{tikzcd}[column sep=small, row sep=small, nodes={font=\scriptsize}]
& (X \otimes Y) \otimes Z \arrow[rr, "c_{X \otimes Y{,}Z}"] & &
Z \otimes (X \otimes Y) \arrow[dr, "a_{Z{,}X{,}Y}"] & \\
X \otimes (Y \otimes Z) \arrow[ur, "a_{X{,}Y{,}Z}"]
\arrow[dr, "\mathrm{id}_{X} \otimes c_{Y{,}Z}"'] & & & &
(Z \otimes X) \otimes Y \\
& X \otimes (Z \otimes Y) \arrow[rr, "a_{X{,}Z{,}Y}"'] & &
(X \otimes Z) \otimes Y \arrow[ur, "c_{X{,}Z} \otimes \mathrm{id}_{Y}"'] &
\end{tikzcd}

A $\otimes$-category $C$ together with an associativity constraint $a$ and
an unity constraint $(1, g, d)$ is a $\otimes$-AU category if all the
following triangles commute

\begin{tikzcd}[column sep=small, row sep=small, nodes={font=\scriptsize}]
X \otimes (1 \otimes Y) \arrow[rr, "a_{X{,}1{,}Y}"] & &
(X \otimes 1) \otimes Y \\
& X \otimes Y \arrow[ul, "\mathrm{id}_{X} \otimes g_{Y}"]
\arrow[ur, "d_{X} \otimes \mathrm{id}_{Y}"'] &
\end{tikzcd}

\page{48}
A $\otimes$-ACU category is a $\otimes$-AC and AU category. An object $X$
of a $\otimes$-ACU category $C$ is \emph{invertible} if there are two
objects $X', X'' \in \mathrm{Ob}\,C$ such that
\[
X' \otimes X \simeq X \otimes X'' \simeq 1.
\]

A \emph{$\otimes$-functor} from a $\otimes$-category $C$ to a
$\otimes$-category $C'$ is a pair $(F, \check{F})$ where $F$ is a functor
$C \to C'$ and $\check{F}$ an isomorphism of bifunctors
\[
\check{F}_{X,Y} : FX \otimes FY \longrightarrow F(X \otimes Y), \qquad
X, Y \in \mathrm{Ob}\,C.
\]

A $\otimes$-functor $(F, \check{F})$ from a $\otimes$-associative category
$C$ to a $\otimes$-associative category $C'$ is \emph{associative} if the
following diagram commutes:

\begin{tikzcd}[column sep=small, row sep=small, nodes={font=\scriptsize}]
FX \otimes (FY \otimes FZ) \arrow[r, "\mathrm{id} \otimes \check{F}"]
\arrow[d, "a'"'] & FX \otimes F(Y \otimes Z) \arrow[r, "\check{F}"] &
F(X \otimes (Y \otimes Z)) \arrow[d, "Fa"] \\
(FX \otimes FY) \otimes FZ \arrow[r, "\check{F} \otimes \mathrm{id}"'] &
F(X \otimes Y) \otimes FZ \arrow[r, "\check{F}"'] &
F((X \otimes Y) \otimes Z)
\end{tikzcd}

where $a$ is the associativity constraint of $C$ and $a'$ of $C'$.

A $\otimes$-functor $(F, \check{F})$ from a $\otimes$-commutative category
$C$ to a $\otimes$-commutative category $C'$ is \emph{commutative} if the
following diagram commutes

\begin{tikzcd}[column sep=normal, row sep=small, nodes={font=\scriptsize}]
FX \otimes FY \arrow[r, "\check{F}"] \arrow[d, "c'"'] &
F(X \otimes Y) \arrow[d, "Fc"] \\
FY \otimes FX \arrow[r, "\check{F}"'] & F(Y \otimes X)
\end{tikzcd}

$c$ and $c'$ being the commutativity constraints of $C$ and $C'$
respectively.

A $\otimes$-functor $(F, \check{F})$ from a $\otimes$-category $C$ with an
unity constraint $(1, g, d)$ to a $\otimes$-category $C'$ with an unity
constraint $(1', g', d')$ is a \emph{$\otimes$-unifer functor} if there
exists an isomorphism $\hat{F} : 1' \simeq F1$ such that the following
diagrams commute:

\begin{tikzcd}[column sep=normal, row sep=small, nodes={font=\scriptsize}]
1' \otimes FX \arrow[r, "\hat{F} \otimes \mathrm{id}_{FX}"] &
F1 \otimes FX \arrow[d, "\check{F}"] \\
FX \arrow[u, "g'_{FX}"] \arrow[r, "Fg_{X}"'] & F(1 \otimes X)
\end{tikzcd}

\begin{tikzcd}[column sep=normal, row sep=small, nodes={font=\scriptsize}]
FX \otimes 1' \arrow[r, "\mathrm{id}_{FX} \otimes \hat{F}"] &
FX \otimes F1 \arrow[d, "\check{F}"] \\
FX \arrow[u, "d'_{FX}"] \arrow[r, "Fd_{X}"'] & F(X \otimes 1)
\end{tikzcd}

\page{49}
It follows from the definition that the isomorphism $\hat{F} : 1' \simeq
F1$, if it exists, is unique.

A $\otimes$-AC functor is a $\otimes$-associative and commutative functor.
A $\otimes$-ACU functor is a $\otimes$-associative, commutative and
unifer functor.

Let $(F, \check{F})$ and $(G, \check{G})$ be $\otimes$-functors from a
$\otimes$-category $C$ to a $\otimes$-category $C'$. A
\emph{$\otimes$-morphism} from the $\otimes$-functor $(F, \check{F})$ to
the $\otimes$-functor $(G, \check{G})$ is a morphism of functors $\lambda
: F \to G$ such that the following diagram commutes

\begin{tikzcd}[column sep=normal, row sep=small, nodes={font=\scriptsize}]
FX \otimes FY \arrow[r, "\check{F}"]
\arrow[d, "\lambda_{X} \otimes \lambda_{Y}"'] &
F(X \otimes Y) \arrow[d, "\lambda_{X \otimes Y}"] \\
GX \otimes GY \arrow[r, "\check{G}"'] & G(X \otimes Y)
\end{tikzcd}

$X, Y \in \mathrm{Ob}\,C$.

Chapter II is a study of Gr-categories and Pic-categories. A
\emph{Gr-category} is a $\otimes$-AU category, the objects of which are
all invertible, and the base category a groupoid (i.e all arrows are
isomorphisms). Thus a Gr-category is like a group. We obtain from this
definition that if $P$ is a Gr-category, the set $\Pi_{0}(P)$ of the
classes up to isomorphism of objects of $P$, together with the operation
induced by the law $\otimes$ of $P$, is a group ; the group
$\mathrm{Aut}(1) = \Pi_{1}(P)$ is a commutative group ; and for all $X \in
\mathrm{Ob}\,P$
\[
\gamma_{X} : u \longmapsto u \otimes \mathrm{id}_{X} \ = \
\mathrm{Aut}(1) \xrightarrow{\ \sim\ } \mathrm{Aut}(X),
\]
\[
\delta_{X} : u \longmapsto \mathrm{id}_{X} \otimes u \ = \
\mathrm{Aut}(1) \xrightarrow{\ \sim\ } \mathrm{Aut}(X).
\]
We attribute thus to a Gr-category $P$ two groups $\Pi_{0}(P)$ and
$\Pi_{1}(P)$ where $\Pi_{1}(P)$ is commutative. Furthermore we can define
an action of $\Pi_{0}(P)$ on $\Pi_{1}(P)$ by the formula
\[
su = \delta_{X}^{-1} \gamma_{X}(u)
\]
for $s \in \Pi_{0}(P)$ represented by $X$ and $u \in \Pi_{1}(P)$. The
commutative group $\Pi_{1}(P)$ together with this action is a left
$\Pi_{0}(P)$-module.

Let $M$ be a group, $N$ a left $M$-module. A \uncertain{preepinglage} of
\note{le mot souligné se lit « preepinglages », avec un s ; le « of » qui
le suit est répété au tournant de page}

\page{50}
type $(M,N)$ for a Gr-category $P$ is a pair $\varepsilon = (\varepsilon_{0},
\varepsilon_{1})$ of isomorphisms
\[
\varepsilon_{0} : M \xrightarrow{\ \sim\ } \Pi_{0}(P), \qquad
\varepsilon_{1} : N \xrightarrow{\ \sim\ } \Pi_{1}(P)
\]
compatible with the action of $M$ on $N$, $\Pi_{0}(P)$ on $\Pi_{1}(P)$. A
Gr-category \emph{preepingled} of type $(M,N)$ is a Gr-category $P$
together with a preepinglage. Finally, an \emph{arrow} of Gr-categories
preepingled of type $(M,N)$, $(P, \varepsilon) \to (P', \varepsilon')$, is
a $\otimes$-associative functor such that the following triangles
commute:

\begin{tikzcd}[column sep=small, row sep=small, nodes={font=\scriptsize}]
\Pi_{0}(P) \arrow[rr] & & \Pi_{0}(P') \\
& M \arrow[ul, "\varepsilon_{0}"] \arrow[ur, "\varepsilon'_{0}"'] &
\end{tikzcd}

\begin{tikzcd}[column sep=small, row sep=small, nodes={font=\scriptsize}]
\Pi_{1}(P) \arrow[rr] & & \Pi_{1}(P') \\
& N \arrow[ul, "\varepsilon_{1}"] \arrow[ur, "\varepsilon'_{1}"'] &
\end{tikzcd}

It follows from this definition that a such arrow is a
$\otimes$-equivalence. Thus the set of the equivalence classes of
Gr-categories preepingled of type $(M,N)$ is equal to the set of the
connected components of the category of Gr-categories preepingled of type
$(M,N)$.

If we consider the cohomology group $H^{3}(M,N)$ of the group $M$ with
coefficients $N$ (in the sense of the group \uncertain{cohomoly} [12]) we
obtain a canonical bijection between the set $H^{3}(M,N)$ and the set of
the equivalence classes of Gr-categories preepingled of type $(M,N)$.
\note{sic pour « cohomoly »}

A \emph{Pic-category} is a Gr-category together with a commutativity
constraint which is compatible with its associativity constraint, i.e the
hexagon axiom is satisfied. Thus a Pic-category is like a commutative
group. We verify immediately that a necessary condition for the existence
of a Pic-category structure on a Gr-category is that $\Pi_{0}(P)$ must be
commutative and act trivially on $\Pi_{1}(P)$. A Pic-category is
\emph{strict} if its commutativity constraint is strict.

Let $M$, $N$ be abelian groups. A preepinglage of type $(M,N)$ for a
Pic-category $P$ is a pair $\varepsilon = (\varepsilon_{0},
\varepsilon_{1})$ of isomorphisms
\[
\varepsilon_{0} : M \xrightarrow{\ \sim\ } \Pi_{0}(P), \qquad
\varepsilon_{1} : N \xrightarrow{\ \sim\ } \Pi_{1}(P).
\]

\page{51}
A Pic-category preepingled of type $(M,N)$ is a Pic-category together
with a preepinglage. We define the \emph{arrow} of such objects in the
same way as for Gr-categories.
\note{le « preepingled » de la première ligne est souligné d'un trait
ondulé — une marque de lecteur}

For next propositions, let us consider two complexes of free abelian
groups
\[
L_{.}(M) : L_{3}(M) \xrightarrow{\ d_{3}\ } L_{2}(M)
\xrightarrow{\ d_{2}\ } L_{1}(M) \xrightarrow{\ d_{1}\ } L_{0}(M)
\longrightarrow M
\]
\[
{}'L_{.}(M) : {}'L_{3}(M) \xrightarrow{\ 'd_{3}\ } {}'L_{2}(M)
\xrightarrow{\ 'd_{2}\ } {}'L_{1}(M) \xrightarrow{\ 'd_{1}\ } {}'L_{0}(M)
\longrightarrow M
\]
where
\[
L_{0}(M) = {}'L_{0}(M) = \mathbf{Z}[M]
\]
\[
L_{1}(M) = {}'L_{1}(M) = \mathbf{Z}[M \times M]
\]
\[
L_{2}(M) = {}'L_{2}(M) = \mathbf{Z}[M \times M \times M] +
\mathbf{Z}[M \times M]
\]
\[
L_{3}(M) = {}'L_{3}(M) + \mathbf{Z}[M]
\]
\[
{}'L_{3}(M) = \mathbf{Z}[M \times M \times M \times M] +
\mathbf{Z}[M \times M \times M] + \mathbf{Z}[M \times M]
\]
\[
d_{1}[x,y] = {}'d_{1}[x,y] = [y] - [x+y] + [x]
\]
\[
d_{2}[x,y] = {}'d_{2}[x,y] = [x,y] - [y,x]
\]
\[
d_{2}[x,y,z] = {}'d_{2}[x,y,z] = [y,z] - [x+y,z] + [x,y+z] - [x,y]
\]
\[
d_{3}[x,y,z,t] = {}'d_{3}[x,y,z,t] = [y,z,t] - [x+y,z,t] + [x,y+z,t] -
[x,y,z+t] + [x,y,z]
\]
\[
d_{3}[x,y,z] = {}'d_{3}[x,y,z] = [x,y,z] - [x,z,y] + [z,x,y] - [y,z] +
[x+y,z] - [x,z]
\]
\[
d_{3}[x,y] = [x,y] + [y,x] = {}'d_{3}[x,y]
\]
\[
d_{3}[x] = [x,x],
\]
so that $L_{.}(M)$ is a truncated resolution of $M$. One obtains a
canonical bijection between the set of the equivalence classes of

\page{52}
Pic-categories preepingled of type $(M,N)$ and the set
$H^{2}(\mathrm{Hom}({}'L_{.}(M), N))$. The exactitude of the complex
$L(M)$ gives us the triviality of the classification of Pic-categories
preepingled of type $(M,N)$ which are strict, i.e all Pic-categories
preepingled of type $(M,N)$ which are strict, are equivalent.
\note{en tête de page, une ligne au crayon à demi rognée par le bord — on
y lit « déterminer $H^{2}(\mathrm{Hom}({}'L_{.}(M), N))$ », la suite
échappe}

Finally chapter III gives us the construction of the solution of two
universal problems : \emph{problem of making objects « unity objects »}
and \emph{problem of reversing objects}.

Let $A$ be a $\otimes$-AC category, $A'$ another $\otimes$-AC category
whose base category is a groupoid, and $(T, \check{T}) : A' \to A$ a
$\otimes$-AC functor. We try to make the objects $TA'$ of $A$, $A' \in
\mathrm{Ob}\,A'$, « unity object », i.e we try to get :
\note{en marge, deux mots au crayon — « auxiliary units » ? — d'une
lecture incertaine}

1° A $\otimes$-ACU category $P$

2° A $\otimes$-AC functor $(D, \check{D}) : A \to P$

3° A $\otimes$-isomorphism
\[
\lambda : (D, \check{D}) \circ (T, \check{T}) \xrightarrow{\ \sim\ }
(I_{P}, \check{I}_{P}),
\]
where $(I_{P}, \check{I}_{P})$ is the $\otimes$-constant functor $1_{P}$
from $A'$ to $P$. The triple $(P, (D, \check{D}), \lambda)$ must be
universal for triples $(Q, (E, \check{E}), \mu)$ satisfying 1°, 2°, 3°.

For the description of the triple $(P, (D, \check{D}), \lambda)$, we
introduce a quotient category of a $\otimes$-AC category as follows :

Let $A$ be a $\otimes$-AC category, $S$ a \emph{multiplicative subset} of
$A$ (that means a subset of the set of all \struck{endomorphisms} arrows
of $A$ such that $\mathrm{id}_{X} \in S$ for all $X \in \mathrm{Ob}\,A$
and the tensor product of two arrows of $S$ belongs to $S$). \note{«
arrows » est écrit au crayon au-dessus d'« endomorphisms » barré ; plus
bas en marge, deux mots au crayon — « what about composition ? » — d'une
lecture incertaine} The $\otimes$-AC category quotient $A^{S}$ of $A$
with respect to $S$ is the solution of the universal problem
\[
(K, \check{K}) : A \longrightarrow B, \qquad
K(u) = \mathrm{id} \ \text{for all} \ u \in S
\]
where $B$ is a $\otimes$-AC category and $(K, \check{K})$ a $\otimes$-AC
functor.

\page{53}
Now let us give an idea of the construction of the triple $(P, (D,
\check{D}), \lambda)$ for $A' \neq \emptyset$ :

1° $\mathrm{Ob}\,P = \mathrm{Ob}\,A$

2° $\mathrm{Hom}_{P}(A,B) = \Phi(A,B)/R_{A,B}$, \quad $A, B \in
\mathrm{Ob}\,P$,

$\Phi(A,B)$ being the set of all triples $(A', B', u)$ where $A', B' \in
\mathrm{Ob}\,A'$, $u \in \mathrm{Fl}\,A$, $u : A \otimes TA' \to B \otimes
TB'$ ; $R_{A,B}$ the equivalence relation defined in $\Phi(A,B)$ as
follows
\[
(A'_{1}, B'_{1}, u) \ R_{A,B} \ (A'_{2}, B'_{2}, u)
\]
\note{le troisième terme s'écrit $u$ dans les deux triples ; le diagramme
les distingue en $u_{1}$ et $u_{2}$}
if and only if there are objects $C'_{1}$, $C'_{2}$ and isomorphisms
\[
u' : A'_{1} \otimes C'_{1} \xrightarrow{\ \sim\ } A'_{2} \otimes C'_{2},
\qquad
v' : B'_{1} \otimes C'_{1} \xrightarrow{\ \sim\ } B'_{2} \otimes C'_{2}
\]
of $A'$ such that the following diagram commutes in $A^{S}$, $\otimes$-AC
quotient category of $A$ with respect to the multiplicative subset of $A$
generated by the endomorphisms of the form $T(c_{A',A'})$ :

\begin{tikzcd}[column sep=small, row sep=small, nodes={font=\scriptsize}]
A \otimes T(A'_{1} \otimes C'_{1})
\arrow[r, "\mathrm{id} \otimes \check{T}^{-1}"]
\arrow[d, "\mathrm{id} \otimes Tu'"'] &
A \otimes (TA'_{1} \otimes TC'_{1}) \arrow[r, "a"] &
(A \otimes TA'_{1}) \otimes TC'_{1} \arrow[r, "u_{1} \otimes \mathrm{id}"]
& (B \otimes TB'_{1}) \otimes TC'_{1} \arrow[d, "a^{-1}"] \\
A \otimes T(A'_{2} \otimes C'_{2})
\arrow[d, "\mathrm{id} \otimes \check{T}^{-1}"'] & & &
B \otimes (TB'_{1} \otimes TC'_{1}) \arrow[d, "\mathrm{id} \otimes \check{T}"] \\
A \otimes (TA'_{2} \otimes TC'_{2}) \arrow[d, "a"'] & & &
B \otimes T(B'_{1} \otimes C'_{1}) \arrow[d, "\mathrm{id} \otimes Tv'"] \\
(A \otimes TA'_{2}) \otimes TC'_{2}
\arrow[r, "u_{2} \otimes \mathrm{id}"'] &
(B \otimes TB'_{2}) \otimes TC'_{2} \arrow[r, "a^{-1}"'] &
B \otimes (TB'_{2} \otimes TC'_{2})
\arrow[r, "\mathrm{id} \otimes \check{T}^{-1}"'] &
B \otimes T(B'_{2} \otimes C'_{2})
\end{tikzcd}

\note{le dernier isomorphisme du bas se lit $\mathrm{id} \otimes
\check{T}^{-1}$ sur la page, là où l'on attendrait $\mathrm{id} \otimes
\check{T}$ ; le trait est trop court pour trancher}

We denote by $[A', B', u]$ the class which has $(A', B', u)$ as
representative.

3° Composition of arrows in $P$. Let $[A', B', u] : A \to B$, $[B'', C'',
v] : B \to C$ be arrows in $P$. We define
\[
[B'', C'', v] \circ [A', B', u] = [A' \otimes B'',\ B' \otimes C'',\
\omega] : A \longrightarrow C
\]

\page{54}
where $\omega$ is such that the following diagram commutes :

\begin{tikzcd}[column sep=small, row sep=small, nodes={font=\scriptsize}]
A \otimes T(A' \otimes B'')
\arrow[r, "\mathrm{id} \otimes \check{T}^{-1}"]
\arrow[dddd, "\omega"'] &
A \otimes (TA' \otimes TB'') \arrow[r, "a"] &
(A \otimes TA') \otimes TB'' \arrow[r, "u \otimes \mathrm{id}"] &
(B \otimes TB') \otimes TB'' \arrow[d, "a^{-1}"] \\
& & & B \otimes (TB' \otimes TB'') \arrow[d, "\mathrm{id} \otimes c"] \\
& & & B \otimes (TB'' \otimes TB') \arrow[d, "a"] \\
& & & (B \otimes TB'') \otimes TB' \arrow[d, "v \otimes \mathrm{id}"] \\
C \otimes T(B' \otimes C'') &
C \otimes (TB' \otimes TC'') \arrow[l, "\mathrm{id} \otimes \check{T}"] &
C \otimes (TC'' \otimes TB') \arrow[l, "\mathrm{id} \otimes c"] &
(C \otimes TC'') \otimes TB' \arrow[l, "a^{-1}"]
\end{tikzcd}

4° $\otimes$-Structure on $P$
\[
A \otimes E \ (\text{in } P) = A \otimes E \ (\text{in } A)
\]
\[
[A', B', u] \otimes [E', F', v] = [A' \otimes E',\ B' \otimes F',\ W]
\]
where $W$ is defined by the commutative diagram (1).

5° ACU constraint in $P$
\[
\bigl([A', A', a \otimes \mathrm{id}],\ [A', A', c \otimes \mathrm{id}],\
(1_{P} = TA'_{0},\ g_{A} = [A'_{0} \otimes A', A', t_{A}],\
d_{A} = [A'_{0} \otimes A', A', r_{A}])\bigr)
\]
where $A'_{0}$ is a fixed object of $A'$, $A'$ an arbitrary object of
$A'$, $g_{A}$ and $d_{A}$ natural isomorphisms
\[
g_{A} : A \longrightarrow 1_{P} \otimes A, \qquad
d_{A} : A \longrightarrow A \otimes 1_{P}
\]
with $t_{A}$ and $r_{A}$ defined by the commutative diagrams (2).
\note{la dernière lettre du crochet de $d_{A}$ est rognée par la marge ;
nous la restituons d'après le diagramme (2)}

6° $(D, \check{D})$ is defined by
\[
DA = A, \qquad Du = [A', A', u \otimes \mathrm{id}_{TA'}], \qquad
\check{D}_{A,B} = \mathrm{id}_{A \otimes B}
\]

\page{55}
\note{les deux diagrammes portent en marge les numéros (1) et (2) que le
texte des pages précédentes appelle}

\begin{tikzcd}[column sep=small, row sep=small, nodes={font=\scriptsize}]
(A \otimes TA') \otimes (E \otimes TE') \arrow[r, "u \otimes v"]
\arrow[d, "a"'] &
(B \otimes TB') \otimes (F \otimes TF') \arrow[d, "a"] \\
((A \otimes TA') \otimes E) \otimes TE'
\arrow[d, "a^{-1} \otimes \mathrm{id}"'] &
((B \otimes TB') \otimes F) \otimes TF'
\arrow[d, "a^{-1} \otimes \mathrm{id}"] \\
(A \otimes (TA' \otimes E)) \otimes TE'
\arrow[d, "(\mathrm{id} \otimes c) \otimes \mathrm{id}"'] &
(B \otimes (TB' \otimes F)) \otimes TF'
\arrow[d, "(\mathrm{id} \otimes c) \otimes \mathrm{id}"] \\
(A \otimes (E \otimes TA')) \otimes TE'
\arrow[d, "a \otimes \mathrm{id}"'] &
(B \otimes (F \otimes TB')) \otimes TF'
\arrow[d, "a \otimes \mathrm{id}"] \\
((A \otimes E) \otimes TA') \otimes TE' \arrow[d, "a^{-1}"'] &
((B \otimes F) \otimes TB') \otimes TF' \arrow[d, "a^{-1}"] \\
(A \otimes E) \otimes (TA' \otimes TE')
\arrow[d, "\mathrm{id} \otimes \check{T}"'] &
(B \otimes F) \otimes (TB' \otimes TF')
\arrow[d, "\mathrm{id} \otimes \check{T}"] \\
(A \otimes E) \otimes T(A' \otimes E') \arrow[r, "W"'] &
(B \otimes F) \otimes T(B' \otimes F')
\end{tikzcd}

\begin{tikzcd}[column sep=small, row sep=small, nodes={font=\scriptsize}]
A \otimes (TA'_{0} \otimes TA') \arrow[r, "\mathrm{id} \otimes \check{T}"]
\arrow[d, "a"'] & A \otimes T(A'_{0} \otimes A') \arrow[d, "t_{A}"] \\
(A \otimes TA'_{0}) \otimes TA' \arrow[r, "c \otimes \mathrm{id}"'] &
(TA'_{0} \otimes A) \otimes TA'
\end{tikzcd}

\begin{tikzcd}[column sep=small, row sep=small, nodes={font=\scriptsize}]
A \otimes (TA'_{0} \otimes TA') \arrow[r, "\mathrm{id} \otimes \check{T}"]
\arrow[d, "a"'] & A \otimes T(A'_{0} \otimes A') \arrow[d, "r_{A}"] \\
(A \otimes TA'_{0}) \otimes TA' \arrow[r, no head] &
(A \otimes TA'_{0}) \otimes TA'
\end{tikzcd}

\note{dans le second carré, le trait horizontal du bas est un double trait
d'égalité sur la page}

7° The $\otimes$-isomorphism
\[
\lambda : (D, \check{D}) \circ (T, \check{T}) \xrightarrow{\ \sim\ }
(I_{P}, \check{I}_{P})
\]
is defined by natural isomorphisms
\[
DTA' = TA'
\xrightarrow{\ \lambda_{A'} = [A'_{0},\, A',\, c_{TA', TA'_{0}}]\ }
I_{P} A' = TA'_{0}, \qquad A' \in \mathrm{Ob}\,A'
\]

$P$ is called the $\otimes$-ACU category of the $\otimes$-AC category $A$
with respect to $(A', (T, \check{T}))$.
\note{un point d'interrogation au crayon précède cette phrase, « ⊗-ACU
category » y est cerclé, et « examples ? » est jeté à l'encre à la suite}

For the problem of reversing objects, let us consider a $\otimes$-category

\page{56}
$C$ with a ACU constraint $(a, c, (1, g, d))$, a $\otimes$-category $C'$
with a ACU constraint $(a', c', (1', g', d'))$, the base category of which
is a groupoid, and a $\otimes$-ACU functor $(F, \check{F}) : C' \to C$. We
try to find a $\otimes$-ACU category $P$ and a $\otimes$-ACU functor $(Q,
\check{Q}) : C \to P$ having the following properties :

1° $QFX'$ is invertible in $P$ for all $X' \in \mathrm{Ob}\,C'$.

2° For all $\otimes$-ACU functor $(Y, \check{Y})$ from $C$ to a
$\otimes$-ACU category $Q$ such that $YFX'$ is invertible in $Q$ for all
$X' \in \mathrm{Ob}\,C'$, there exists a $\otimes$-ACU functor $(E,
\check{E})$, unique up to $\otimes$-isomorphism, from $P$ to $Q$ such
that $(Y, \check{Y}) \simeq (E, \check{E}) \circ (Q, \check{Q})$.
\note{sur la page, les catégories sont soulignées : le foncteur $(Q,
\check{Q})$ et la catégorie $Q$ du 2° se distinguent ainsi. En marge,
trois mots au crayon dont seuls « équivalence » et « catégories » se
laissent lire}

This problem is reduced by the first by putting $A' = C'$, $A = C \times
C'$, $TX' = (FX', X')$ and by remarking that if $C$, $C'$, $Q$ are
$\otimes$-ACU categories, $\mathrm{Hom}^{\otimes \mathrm{ACU}}(C, Q)$ the
category of all $\otimes$-ACU functor from $C$ to $Q$, then there is a
canonical equivalence of categories
\[
\mathrm{Hom}^{\otimes \mathrm{ACU}}(C \times C', Q) \longrightarrow
\mathrm{Hom}^{\otimes \mathrm{ACU}}(C, Q) \times
\mathrm{Hom}^{\otimes \mathrm{ACU}}(C', Q).
\]
The $\otimes$-ACU category $P$ thus defined is called the
\emph{$\otimes$-category of fractions} of the category $C$ with respect
to $(C', (F, \check{F}))$. The $\otimes$-category of fractions of $C$
with respect to $(C^{is}, (\mathrm{id}_{C^{is}}, \mathrm{id}))$ is a
Pic-category which is called the \emph{Pic-envelope} of the category $C$,
and denoted by $\mathrm{Pic}(C)$.

For an application of the Pic-envelope, we take $C = P(R)$, category of
all finitely generated $R$-modules ($R$ a ring) and $P =
\mathrm{Pic}(P(R))$, then one obtains
\[
\Pi_{0}(P) \simeq K^{0}(R), \qquad \Pi_{1}(P) \simeq K^{1}(R)
\]
where $K^{0}(R)$ is the Grothendieck group and $K^{1}(R)$ the Whitehead
group [1].
\marginal{this Pic-envelope is not a strict Pic-category}
\note{la note marginale, en anglais, est de la main de l'annotateur}

The use of the $\otimes$-category of fractions of a $\otimes$-ACU
category gives us the following result :

\page{57}
Let $C$ be a $\otimes$-ACU category, $Z$ an arbitrary object of $C$
different from the unity object $1$, $S$ the functor from $C$ to $C$
defined by $X \mapsto X \otimes Z$. The \emph{suspension category} of the
$\otimes$-ACU category $C$ defined by the object $Z$ is the triple $(P,
i, p)$ which solves the universal problem for triples $(Q, j, q)$, where
$Q$ is a category, $j$ a functor from $C$ to $Q$, and $q$ an equivalence
of categories from $Q$ to $Q$, so that the following diagram commutes

\begin{tikzcd}[column sep=normal, row sep=small, nodes={font=\scriptsize}]
C \arrow[r, "S"] \arrow[d, "j"'] & C \arrow[d, "j"] \\
Q \arrow[r, "q"'] & Q
\end{tikzcd}

up to natural isomorphism. In the case where $C$ is the homotopy category
of pointed topological spaces $\mathrm{Htp}_{*}$ together with the smash
$\wedge$ (the smash $X \wedge Y$ of two spaces $X$ and $Y$, with the base
points $x_{0}$ and $y_{0}$, is obtained from the product $X \times Y$ by
shrinking the subset $X \vee Y = X \times \{y_{0}\} \cup \{x_{0}\} \times
Y$ to a single point which is taken as the base point of $X \wedge Y$),
and the usual ACU constraint ; and $Z$ is the 1-sphere $S^{1}$ hence $S$
is the suspension functor, we get the well-known definition of the
suspension category.

Let $C'$ be the $\otimes$-stable subcategory of $C$ generated by $Z$ and
$P$ the $\otimes$-category of fractions of $C$ with respect to $(C', (F,
\mathrm{id}))$ where $F : C' \to C$ is the inclusion functor. One obtains
a functor $G : P \to P$ from the suspension category to the
$\otimes$-category of fractions $P$. \note{les deux $P$ — la catégorie de
suspension et la catégorie de fractions — s'écrivent de la même lettre sur
la page ; la phrase les distingue} If $G$ is not faithful, that is the
case of the homotopy category of pointed topological spaces
$\mathrm{Htp}_{*}$ together with the smash $\wedge$ and the 1-sphere
$S^{1}$ ; then it is impossible to construct in $P$ a law $\otimes$ such
that $P$ together with this law is a $\otimes$-ACU category, $iZ$
invertible in $P$, and $i$ imbedded in a pair $(i, \check{\imath})$ which
is a $\otimes$-ACU functor from $C$ to $P$.

\page{58}
\textbf{References}

[1] Bass, H.\ : K-theory and stable algebra. Publ.\ math.\ de l'IHES,
n° 22.

[2] Bénabou, J.\ : Thèse, Paris 1966.

[3] Bourbaki : Théorie des ensembles.

[4] --- : Algèbre commutative.

[5] --- : Algèbre multilinéaire.

[6] Deligne, P.\ : Champs de Picard strictement commutatifs. SGA 4 XVIII.

[7] Eilenberg, S.\ and Kelly, G.M : Closed categories. Proceedings of the
conference on categorical algebra (421-561). Springer-Verlag 1965.

[8] Freyd, P.\ : Stable homotopy. Proceedings of the conference on
categorical algebra (121-176). Springer-Verlag 1965.

[9] Grothendieck, A.\ : Biextensions de faisceaux de groupes. SGA 7 VII.

[10] --- : Catégories cofibrées additives et complexe cotangent relatif.
Lecture notes in mathematics N° 79, Springer-Verlag 1968.

[11] Mac-Lane, S.\ : Homology. Springer-Verlag 1967.

[12] --- : Categorical algebra. Bull.\ Amer.\ Mat.\ Soc.\ 71 (1965)
(40-106).

[13] Mitchell, B.\ : Theory of categories. Academic Press 1965.

[14] Neantro Saavedra Rivano : Thèse, Paris (1970 ?)

[15] --- : Catégories \uncertain{tanakiennes}. Lecture notes in
mathematics N° 265. Springer-Verlag 1972. \note{sic : un seul n}

[16] Spanier, E : Algebraic topology. Mc Graw-Hill Inc.\ 1966.

\end{document}
